\chapter{BERTAHAN TANPA IDENTITAS}

\section{Unitalisasi Aljabar Banach}
\begin{defn}\label{00150005} Misalkan $A$ suatu aljabar. Definisikan
 \index{<binopunit@$A \bowtie \C$ (unitalisasi suatu aljabar)}%
$A \bowtie \C$ sebagai himpunan $A \times \C$ yang penjumlahan dan perkalian skalarnya
didefinisikan secara titik demi titik dan perkaliannya didefinisikan oleh
  \[ (a,\lambda)\cdot(b,\mu) = (ab + \mu a + \lambda b, \lambda\mu). \]
Jika aljabar $A$ dilengkapi dengan suatu involusi, definisikan involusi pada $A \bowtie \C$
secara titik demi titik; yaitu, $(a,\lambda)^* := (a^*, \conj\lambda)$. (Notasi $A \bowtie \C$
tidak baku.)
\end{defn}

\begin{prop}\label{0015001} Jika $A$ suatu aljabar, maka $A \bowtie \C$ adalah aljabar
beridentitas yang di dalamnya $A$ dibenamkan sebagai (yaitu, isomorfik dengan) suatu ideal
sedemikian sehingga $(A \bowtie \C)/A \cong \C$. Identitas $A \bowtie \C$ adalah $(0,1)$. Jika
$A$ suatu aljabar-$*\,$, maka $A \bowtie \C$ adalah aljabar-$*\,$ beridentitas yang di dalamnya
$A$ merupakan ideal-$*\,$ sedemikian sehingga $(A \bowtie \C)/A \cong \C$.
\end{prop}

\begin{defn} Aljabar $A \bowtie \C$ adalah
 \index{unitalisasi!suatu aljabar-$*\,$}%
\df{unitalisasi} dari aljabar (atau aljabar-$*\,$)~$A$.
\end{defn}

\begin{notn}\label{00150012} Dalam konstruksi sebelumnya, unsur-unsur aljabar
$A \bowtie \C$ secara teknis adalah pasangan terurut $(a,\lambda)$. Unsur-unsur itu biasanya
ditulis secara berbeda. Misalkan $\iota\colon A \sto A \bowtie \C \colon a \mapsto (a,0)$ dan
$\pi\colon A \bowtie \C \sto \C\colon (a,\lambda) \mapsto \lambda$. Karena $(0,1)$ adalah
identitas dalam $A \bowtie \C$, diperoleh
   \begin{align*} (a,\lambda) &= (a,0) + (\vc 0,\lambda) \\
                              &= \iota(a) + \lambda \vc 1_{A \bowtie \C}
   \end{align*}
Sudah lazim untuk memperlakukan $\iota$ sebagai pemetaan inklusi. Jadi, wajar untuk
menulis $(a,\lambda)$ sebagai $a + \lambda \vc 1_{A \bowtie \C}$ atau cukup sebagai
$a + \lambda \vc 1$. Tampaknya tidak timbul ambiguitas dengan menghilangkan acuan pada
identitas perkalian, sehingga notasi baku untuk pasangan $(a,\lambda)$ adalah $a + \lambda$.
\end{notn}

\begin{defn}\label{00150014} Misalkan $A$ suatu aljabar yang \emph{tak beridentitas}. Definisikan
 \index{spektrum!dalam aljabar tak beridentitas}%
\df{spektrum} suatu unsur $a \in A$ sebagai spektrum $a$ ketika dipandang sebagai unsur
dari~$A \bowtie \C$; yaitu, $\sigma_A(a) := \sigma_{A \bowtie \C}(a)$.
\end{defn}

\begin{defn}\label{00150015} Misalkan $A$ suatu aljabar bernorma (dengan atau tanpa involusi). Pada
unitalisasi $A \bowtie \C$ dari $A$, definisikan $\norm{(a,\lambda)} := \norm a + \abs\lambda$.
\end{defn}

\begin{prop}\label{001500151} Misalkan $A$ suatu aljabar bernorma. Pemetaan
 \index{unitalisasi!suatu aljabar bernorma}%
 \index{unitalisasi!suatu aljabar Banach-$*\,$}%
$(a,\lambda) \mapsto \norm a + \abs\lambda$ yang didefinisikan di atas adalah norma yang
menjadikan $A \bowtie \C$ suatu aljabar bernorma. Jika $A$ suatu aljabar Banach (berturut-turut,
aljabar Banach-$*\,$), maka $A \bowtie \C$ adalah aljabar Banach (berturut-turut, aljabar
Banach-$*\,$). Aljabar Banach (atau aljabar Banach-$*\,$) yang dihasilkan adalah
\df{unitalisasi} dari $A$ dan akan
 \index{<unitization@$A_e$ (unitalisasi suatu aljabar Banach)}%
dinotasikan dengan~$A_e$.
\end{prop}

Dengan definisi \emph{spektrum} yang diperluas ini, banyak fakta terdahulu untuk aljabar
Banach beridentitas tetap berlaku dalam tatanan yang lebih umum. Secara khusus, untuk
rujukan selanjutnya, kita nyatakan kembali butir-butir \ref{000661}, \ref{C073134},
\ref{0006703}, dan~\ref{C073157}.

\begin{prop}\label{001500157} Misalkan $a$ suatu unsur dari aljabar Banach~$A$. Maka spektrum
$a$ kompak dan $\abs\lambda \le \norm a$ untuk setiap $\lambda \in \sigma(a)$.
\end{prop}

\begin{prop}\label{00150016} Spektrum setiap unsur dari suatu aljabar Banach tidak kosong.
\end{prop}

\begin{prop}\label{00150017} Untuk setiap unsur $a$ dari suatu aljabar Banach, $\rho(a) \le
\norm a$ dan $\rho(a^n) = \bigr(\rho(a)\bigr)^n$.
\end{prop}

\begin{thm}[Rumus Radius Spektral]\label{00150018} Jika $a$ suatu unsur dari aljabar Banach,
 \index{spektral!radius!rumus untuk}%
maka
   \[ \rho(a) = \inf\bigl\{\norm{a^n}^{1/n}\colon n \in \N\bigr\}
              = \lim_{n \sto \infty} \norm{a^n}^{1/n}. \]
\end{thm}

\begin{defn} Misalkan $A$ suatu aljabar. Suatu ideal kiri $J$ dalam $A$ adalah
 \index{modular!ideal kiri}%
 \index{ideal!kiri modular}%
\df{ideal kiri modular} jika terdapat suatu unsur $u$ dalam $A$ sedemikian sehingga
$au - a \in J$ untuk setiap $a \in A$. Unsur $u$ demikian disebut
 \index{identitas!kanan!terhadap suatu ideal}%
\df{identitas kanan terhadap $J$}. Demikian pula, suatu ideal kanan $J$ dalam $A$ adalah
 \index{modular!ideal kanan}%
 \index{ideal!kanan modular}%
\df{ideal kanan modular} jika terdapat suatu unsur $v$ dalam $A$ sedemikian sehingga
$va - a \in J$ untuk setiap $a \in A$. Unsur $v$ demikian disebut
 \index{identitas!kiri!terhadap suatu ideal}%
\df{identitas kiri terhadap $J$}. Suatu ideal dua sisi $J$ adalah
 \index{modular!ideal}%
 \index{ideal!modular}%
\df{ideal modular} jika terdapat suatu unsur $e$ yang sekaligus merupakan identitas kiri
dan identitas kanan terhadap~$J$.
\end{defn}

\begin{prop} Suatu ideal $J$ dalam suatu aljabar bersifat modular jika dan hanya jika ideal
tersebut merupakan ideal kiri modular sekaligus ideal kanan modular.
\end{prop}

\begin{proof}[\emph{Petunjuk untuk pembuktian}] Tunjukkan bahwa jika $u$ suatu identitas kanan
terhadap $J$ dan $v$ suatu identitas kiri terhadap $J$, maka $vu$ sekaligus merupakan
identitas kanan dan identitas kiri terhadap~$J$. \ns
\end{proof}

\begin{prop} Suatu ideal $J$ dalam suatu aljabar $A$ bersifat modular jika dan hanya jika
aljabar hasil bagi $A/J$ beridentitas.
\end{prop}

\begin{exam} Misalkan $X$ suatu ruang Hausdorff kompak lokal. Untuk setiap $x \in X$, ideal
$J_x$ adalah ideal modular maksimal dalam aljabar-$C^*$ $\fml C_0(X)$ dari fungsi-fungsi
kontinu bernilai kompleks pada~$X$.
\end{exam}

\begin{proof} Menurut versi \emph{lema Urysohn} untuk ruang Hausdorff kompak lokal
(lihat, misalnya, \cite{Erdman:2007}, teorema 17.2.10), terdapat suatu fungsi $g \in
\fml C_0(X)$ sedemikian sehingga $g(x) = 1$. Jadi, $J_x$ bersifat modular karena $g$
merupakan suatu identitas terhadap~$J_x$. Karena $\fml C_0(X) = J_x \oplus\,\, \spn\{g\}$,
ideal $J_x$ mempunyai kodimensi $1$ dan oleh karena itu maksimal.
\end{proof}

\begin{prop} Jika $J$ suatu ideal modular sejati dalam aljabar Banach, maka penutupannya
juga demikian.
\end{prop}

\begin{cor} Setiap ideal modular maksimal dalam suatu aljabar Banach tertutup.
\end{cor}

\begin{prop} Setiap ideal modular sejati dalam suatu aljabar Banach termuat dalam suatu
ideal modular maksimal.
\end{prop}

\begin{prop}\label{00150144} Misalkan $A$ suatu aljabar Banach komutatif dan $\phi \in
\Delta A$. Maka $\ker\phi$ adalah ideal modular maksimal dalam~$A$ dan $A/\ker\phi$ adalah
suatu medan. Lebih lanjut, setiap ideal modular maksimal merupakan kernel dari tepat satu
karakter dalam~$\Delta A$.
\end{prop}

\begin{prop}\label{00150147} Setiap fungsional linear multiplikatif pada suatu aljabar
Banach komutatif $A$ bersifat kontinu. Bahkan, setiap karakter bersifat kontraktif.
\end{prop}

\begin{exam} Misalkan $A_e$ unitalisasi dari suatu aljabar Banach komutatif~$A$ (lihat
proposisi~\ref{001500151}).
 \index{phi@$\phi_\infty$ (suatu karakter pada unitalisasi suatu aljabar)}%
Definisikan
   \[ \phi_\infty\colon A_e \sto \C\colon (a,\lambda) \mapsto \lambda\,. \]
Maka $\phi_\infty$ adalah suatu karakter pada~$A_e$.
\end{exam}

\begin{prop}\label{00150154} Setiap karakter $\phi$ pada suatu aljabar Banach komutatif
$A$ mempunyai perpanjangan tunggal menjadi suatu
 \index{phi@$\phi_e$ (perpanjangan suatu karakter $\phi$ ke unitalisasi suatu aljabar)}%
karakter $\phi_e$ pada unitalisasi $A_e$ dari~$A$. Selain itu, pembatasan pada $A$ dari sebarang
karakter pada $A_e$, dengan pengecualian yang jelas yaitu $\phi_\infty$, adalah suatu
karakter pada~$A$.
\end{prop}

\begin{prop} Jika $A$ suatu aljabar Banach komutatif, maka
 \begin{enumerate}
    \item[(a)] $\Delta A$ adalah ruang Hausdorff kompak lokal,
    \item[(b)] $\Delta A_e = \Delta A \cup \{\phi_\infty\}$,
    \item[(c)] $\Delta A_e$ adalah kompaktifikasi satu titik dari~$\Delta A$, dan
    \item[(d)] pemetaan $\phi \mapsto \phi_e$ adalah homeomorfisme dari $\Delta A$ ke
$\Delta A_e \setminus \{\phi_\infty\}$.
 \end{enumerate}
Jika $A$ beridentitas (sehingga $\Delta A$ kompak), maka $\phi_\infty$ adalah titik
terisolasi dari~$\Delta A_e$.
\end{prop}

\begin{thm} Jika $A$ suatu aljabar Banach komutatif dengan ruang karakter yang tidak kosong,
maka transformasi Gelfand
  \[ \Gamma = \Gamma_A\colon A \sto \fml C_0(\Delta A)\colon a \mapsto \wh a = \Gamma_a \]
adalah suatu homomorfisme aljabar kontraktif dan $\rho(a) = \norm{\wh a}_u$. Lebih lanjut,
jika $A$ \emph{tidak} beridentitas, maka $\sigma(a) = \ran \wh a \cup \{0\}$ untuk setiap
$a \in A$.
\end{thm}

















\section{Barisan Eksak dan Ekstensi}
Prosedur yang cukup sederhana untuk unitalisasi aljabar Banach-$*\,$
(lihat~\ref{00150005} dan~\ref{00150015}) tidak dapat diterapkan begitu saja pada
aljabar-$C^*$. Norma yang didefinisikan dalam~\ref{00150015} tidak memenuhi syarat-$C^*$
(definisi~\ref{0013}). Ternyata unitalisasi aljabar-$C^*$ sedikit lebih rumit. Sebelum
menelaah perinciannya, kita melihat beberapa bahan persiapan mengenai \emph{barisan eksak}
dan \emph{ekstensi} aljabar-$C^*$.

Dalam dunia aljabar-$C^*$, barisan eksak didefinisikan pada dasarnya dengan cara yang sama
seperti barisan eksak aljabar Banach (lihat~\ref{defn_Bsp_exact_seq}).

\begin{defn} Suatu barisan aljabar-$C^*$ dan homomorfisme-$*\,$
   \[\xymatrix{{\cdots}\ar[r] & A_{n-1}\ar[r]^{\phi_n}
          & A_n\ar[r]^{\phi_{n+1}} & A_{n+1}\ar[r] & {\cdots}
   }\]
dikatakan
 \index{barisan!eksak}%
 \index{eksak!barisan}%
\df{eksak di} $A_n$ jika $\ran \phi_n = \ker \phi_{n+1}$. Suatu barisan disebut
\df{eksak} jika barisan tersebut eksak di setiap aljabar-$C^*$ penyusunnya. Suatu barisan
aljabar-$C^*$ dan homomorfisme-$*\,$ berbentuk
   \begin{equation}\label{001500202i2}\xymatrix{
      0 \ar[r] & A \ar[r]^\phi & E\ar[r]^\psi & B\ar[r] & 0
   }\end{equation}
 \index{barisan!eksak pendek}%
 \index{barisan eksak pendek}%
adalah suatu \df{barisan eksak pendek}. (Di sini $\vc 0$ menyatakan aljabar-$C^*$ trivial
berdimensi $0$, dan panah-panah tanpa label adalah homomorfisme-$*$ yang jelas.) Barisan
eksak pendek aljabar-$C^*$~\eqref{001500202i2} adalah
 \index{eksak!barisan!terbelah}%
 \index{barisan!eksak terbelah}%
 \index{eksak terbelah!barisan}%
\df{eksak terbelah} jika terdapat suatu homomorfisme-$*\,$ $\xi\colon B \sto E$ sedemikian
sehingga $\psi\circ\xi = \id B$.
\end{defn}

Definisi-definisi sebelumnya diberikan untuk kategori $\cat{CSA}$ dari aljabar-$C^*$ dan
homomorfisme-$*\,$. Tentu saja, tidak ada yang khusus mengenai kategori tertentu ini.
Barisan eksak bermakna dalam banyak situasi, misalnya dalam berbagai kategori ruang Banach,
aljabar Banach, ruang Hilbert, ruang vektor, grup Abel, modul, dan sebagainya.

Sering kali, dalam konteks aljabar-$C^*$, barisan eksak \eqref{001500202i2} disebut suatu
 \index{ekstensi}%
\df{ekstensi}. Sebagian penulis menyebutnya \emph{ekstensi $A$ oleh $B$} (misalnya,
\cite{Wegge-Olsen:1993} dan \cite{Davidson:1996}), sedangkan penulis lain mengatakan bahwa
barisan itu merupakan \emph{ekstensi $B$ oleh $A$} (\cite{Murphy:1990},
\cite{Fillmore:1996}, dan \cite{Blackadar:2006}). Dalam \cite{Fillmore:1996} dan
\cite{Blackadar:2006}, \emph{ekstensi} didefinisikan sebagai barisan~\eqref{001500202i2};
dalam \cite{Wegge-Olsen:1993}, ekstensi didefinisikan sebagai tripel terurut
$(\phi,E,\psi)$; dan dalam \cite{Murphy:1990} serta \cite{Davidson:1996}, ekstensi
didefinisikan sebagai aljabar-$C^*$ $E$ itu sendiri. Terlepas dari definisi formalnya,
umum dikatakan bahwa \emph{$E$ adalah ekstensi $A$ oleh $B$ (atau $B$ oleh $A$)}.

\begin{conv}\label{00151212} Dalam suatu aljabar-$C^*$, kata \emph{ideal} akan selalu
 \index{konvensi!ideal bersifat swaadjoin}%
 \index{konvensi!ideal bersifat tertutup}%
 \index{ideal!dalam suatu aljabar-$C^*$}%
berarti suatu \hbox{ideal-$*\,$} dua sisi tertutup (kecuali, tentu saja, jika yang
sebaliknya dinyatakan secara eksplisit). Sebentar lagi akan kita tunjukkan (dalam
proposisi~\ref{0019017}) bahwa syarat agar suatu ideal dalam aljabar-$C^*$ bersifat
swaadjoin sebenarnya berlebih. Suatu ideal aljabar dua sisi dari aljabar-$C^*$ yang tidak
harus tertutup akan disebut
 \index{ideal!aljabar}%
 \index{aljabar!ideal}%
\df{ideal aljabar}. Suatu ideal aljabar swaadjoin dari aljabar-$C^*$ akan disebut
 \index{ideal!aljabar-$\ast$}%
 \index{aljabar!ideal-$\ast$}%
\df{ideal-$\ast$ aljabar}.
\end{conv}

Bandingkan dua proposisi berikut dengan~\ref{00151231} dan~\ref{00151232} yang hampir identik.

\begin{prop}\label{00151231c} Tinjau diagram berikut dalam kategori aljabar-$C^*$ dan
homomorfisme-${}^*$
 \[\xymatrix{
      0\ar[r] & A\ar[d]^f\ar[r]^j & B\ar[d]^g\ar[r]^k
              & C\ar@{-->}[d]^h \ar[r] & 0 \\
      0\ar[r] & A'\ar[r]_{j'} & B'\ar[r]_{k'}
      & C'\ar[r] & 0
 }\]
Jika baris-barisnya eksak dan persegi kirinya komutatif, maka terdapat homomorfisme-${}^*$
$h\colon C \sto C'$ tunggal yang membuat persegi kanannya komutatif.
\end{prop}

\begin{prop}[Lema Lima]\label{00151232c}
 \index{lema lima}%
Andaikan dalam diagram aljabar-$C^*$ dan homomorfisme-${}^*$ berikut
 \[\xymatrix{
      0\ar[r] & A\ar[d]^f\ar[r]^j & B\ar[d]^g\ar[r]^k
              & C\ar[d]^h \ar[r] & 0 \\
      0\ar[r] & A'\ar[r]_{j'} & B'\ar[r]_{k'}
      & C'\ar[r] & 0
 }\]
baris-barisnya eksak dan persegi-perseginya komutatif. Buktikan hal-hal berikut.
 \begin{enumerate}
  \item Jika $g$ surjektif, maka $h$ juga surjektif.
  \item Jika $f$ surjektif dan $g$ injektif, maka $h$ injektif.
  \item Jika $f$ dan $h$ surjektif, maka $g$ juga surjektif.
  \item Jika $f$ dan $h$ injektif, maka $g$ juga injektif.
 \end{enumerate}
\end{prop}

\begin{defn} Misalkan $A$ dan $B$ aljabar-$C^*$. Kita mendefinisikan
 \index{jumlah langsung!aljabar-$C^*$}%
 \index{eksternal!jumlah langsung}%
 \index{C*@aljabar-$C^*$!jumlah langsung dari}%
 \index{<binopsum@$A \oplus B$ (jumlah langsung aljabar-$C^*$)}%
\df{jumlah langsung (eksternal)} dari $A$ dan $B$, yang dinotasikan dengan $A \oplus B$,
sebagai hasil kali Cartesius $A \times B$ dengan operasi-operasi aljabar yang didefinisikan
secara titik demi titik dan norma yang diberikan oleh
   \[ \norm{(a,b)} = \max\{\norm a,\norm b\} \]
untuk semua $a \in A$ dan $b \in B$. Notasi alternatif untuk unsur $(a,b)$ dalam
$A \oplus B$
 \index{<binopsum@$a \oplus b$ (suatu unsur dari $A \oplus B$)}%
adalah $a \oplus b$.
\end{defn}

\begin{exam} Misalkan $A$ dan $B$ aljabar-$C^*$. Maka jumlah langsung $A$ dan $B$ adalah
suatu aljabar-$C^*$ dan barisan berikut adalah barisan eksak pendek terbelah:
 \[\xymatrix{
    \vc 0 \ar[r] & A \ar[r]^-{\iota_1} & A \oplus B \ar[r]^-{\pi_2}
             & B \ar[r] & \vc 0
 }\]
\end{exam}
\noindent Pemetaan-pemetaan yang ditunjukkan di atas adalah pemetaan-pemetaan yang jelas:
   \[ \iota_1\colon A \sto A \oplus B\colon a \mapsto (a,\vc 0) \qquad
        \text{dan} \qquad  \pi_2 \colon A \oplus B \sto B: (a,b) \mapsto b\,. \]
 \index{ekstensi!jumlah langsung}%
 \index{jumlah langsung!ekstensi}
Ini adalah \df{ekstensi jumlah langsung}.

\begin{prop} Jika $A$ dan $B$ adalah aljabar-$C^*$ tak nol, maka
 \index{jumlah langsung!aljabar-$C^*$}%
 \index{produk!dalam $\cat{CSA}$}%
 \index{csa@$\cat{CSA}$!produk dalam}%
jumlah langsungnya $A \oplus B$ merupakan produk dalam kategori $\cat{CSA}$ dari
aljabar-$C^*$ dan homomorfisme-$*\,$ (lihat contoh~\ref{0007812}). Jumlah langsung itu
beridentitas jika dan hanya jika $A$ dan $B$ keduanya beridentitas.
\end{prop}

\begin{defn}\label{0015151} Misalkan $A$ dan $B$ aljabar-$C^*$ serta $E$ dan $E'$ merupakan
ekstensi $A$ oleh $B$. Ekstensi-ekstensi ini
 \index{ekuivalensi!kuat!ekstensi}%
 \index{ekstensi!ekuivalensi kuat}%
\df{ekuivalen kuat} jika terdapat suatu isomorfisme-${}^*$ $\theta\colon E \sto E'$
yang membuat diagram
 \[\xymatrix{
     0\ar[r] & A\ar[r]\ar@{=}[d] & E\ar[r]\ar[d]^\theta
             & B\ar[r]\ar@{=}[d] & 0 \\
     0\ar[r] & A\ar[r] & E'\ar[r] & B\ar[r] & 0
 }\]
komutatif.
\end{defn}

\begin{prop}\label{0015152} Dalam definisi sebelumnya, cukup disyaratkan bahwa $\theta$
merupakan suatu homomorfisme-${}^*$.
\end{prop}

\begin{prop}\label{0015155} Misalkan $A$ dan $B$ aljabar-$C^*$. Suatu ekstensi
 \[\xymatrix{
     \vc 0\ar[r] & A\ar[r]^\phi & E\ar[r]^\psi & B\ar[r] & \vc 0
 }\]
ekuivalen kuat dengan ekstensi jumlah langsung $A \oplus B$ jika dan hanya jika terdapat\newline
suatu \mbox{homomorfisme-${}^*$} $\nu\colon E \sto A$ sedemikian sehingga
$\nu\circ\phi = \id A$.
\end{prop}

\begin{prop}\label{0015157} Jika barisan-barisan aljabar-$C^*$
 \begin{equation}\label{0015157i}\xymatrix{
     0\ar[r] & A\ar[r]^\phi & E\ar[r]^\psi & B\ar[r] & 0
 }\end{equation}
dan
 \begin{equation}\label{0015157ii}\xymatrix{
     0\ar[r] & A\ar[r]^{\phi'} & E'\ar[r]^{\psi'} & B\ar[r] & 0
 }\end{equation}
ekuivalen kuat dan \eqref{0015157i} terbelah, maka \eqref{0015157ii} juga terbelah.
\end{prop}

















\section{Unitalisasi Aljabar-$C^*$}
Prosedur yang cukup sederhana untuk unitalisasi aljabar Banach-$*\,$
(lihat~\ref{00150005} dan~\ref{00150015}) tidak dapat diterapkan begitu saja pada
aljabar-$C^*$. Norma yang didefinisikan dalam~\ref{00150015} tidak memenuhi syarat-$C^*$
(definisi~\ref{0013}). Ternyata unitalisasi aljabar-$C^*$ sedikit lebih rumit.

\begin{prop} Kernel dari suatu homomorfisme-$*\,$ $\phi\colon A \sto B$ antara
aljabar-$C^*$ adalah suatu ideal-$*\,$ tertutup dalam~$A$ dan jangkauannya adalah suatu
subaljabar-$*\,$ dari~$B$.
\end{prop}

\begin{prop}\label{001315} Jika $a$ suatu unsur dalam aljabar-$C^*$, maka
  \begin{align*}
    \norm{a} &= \sup\{\norm{xa}\colon \norm{x} \le 1\} \\
             &= \sup\{\norm{ax}\colon \norm{x} \le 1\}.
  \end{align*}
\end{prop}

\begin{cor}\label{0013151} Jika $a$ suatu unsur dari aljabar-$C^*$ $A$, operator $L_a$,
yang disebut
 \index{kiri!operator perkalian}%
 \index{operator!perkalian kiri}%
\df{perkalian kiri oleh~$a$} dan didefinisikan oleh
   \[ L_a\colon A \sto A\colon x \mapsto ax \]
adalah suatu operator (linear terbatas) pada~$A$. Lebih lanjut, pemetaan
   \[ L\colon A \sto \ofml B(A)\colon a \mapsto L_a \]
sekaligus merupakan isometri dan homomorfisme aljabar.
\end{cor}

\begin{prop}\label{0015213} Misalkan $A$ suatu aljabar-$C^*$. Maka terdapat suatu
aljabar-$C^*$ beridentitas $\wt A$ yang di dalamnya $A$ dibenamkan sebagai suatu ideal
sedemikian sehingga barisan
 \begin{equation}\label{0015213i}
   \xymatrix{
     \vc 0\ar[r] & A\ar[r] & \wt A \ar[r] & \C \ar[r] & \vc 0
 }\end{equation}
eksak terbelah. Jika $A$ beridentitas, maka barisan \eqref{0015213i} ekuivalen kuat dengan
ekstensi jumlah langsung, sehingga $\wt A \cong A \oplus \C$. Jika $A$ \emph{tidak}
beridentitas, maka $\wt A$ tidak isomorfik dengan $A \oplus \C$.
\end{prop}

\begin{proof}[\emph{Petunjuk untuk pembuktian}] Pembuktian hasil ini agak rumit. Setiap orang
sebaiknya menelusuri semua perinciannya setidaknya sekali dalam hidupnya. Berikut ini
adalah garis besar suatu pembuktian.

Perhatikan bahwa kita berbicara tentang unitalisasi dari suatu aljabar-$C^*$ $A$
\emph{baik $A$ sudah memiliki unit (identitas perkalian) maupun belum}. Kita membagi
argumen menjadi dua kasus.

\vskip 8 pt

\noindent\textbf{Kasus 1: aljabar $A$ beridentitas.}
\begin{enumerate}
    \item Pada aljabar $A \bowtie \C$, definisikan
       \[ \norm{(a,\lambda)} := \max\{\norm{a + \lambda \vc 1_A}, \abs\lambda\} \]
dan misalkan $\wt A := A \bowtie \C$ bersama dengan norma ini.
    \item Buktikan bahwa pemetaan $(a,\lambda) \mapsto \norm{(a,\lambda)}$ adalah norma
pada~$A \bowtie \C$.
    \item Buktikan bahwa norma ini adalah norma aljabar.
    \item Tunjukkan bahwa norma ini, sesungguhnya, merupakan norma-$C^*$ pada
$A \bowtie \C$.
    \item Amati bahwa norma ini merupakan perpanjangan dari norma pada~$A$.
    \item Buktikan bahwa $A \bowtie \C$ adalah suatu aljabar-$C^*$ dengan memverifikasi
kelengkapan ruang metrik yang diinduksi oleh norma sebelumnya.
    \item Buktikan bahwa barisan
       \[ \vc 0 \to A \to^\iota \wt A \two/->`<-/^Q_\psi \C \to \vc 0  \]
eksak terbelah (dengan $\iota\colon a \mapsto (a,0)$, $Q\colon (a,\lambda) \mapsto
\lambda$, dan $\psi\colon \lambda \mapsto (0,\lambda)$).
    \item Buktikan bahwa $\cstariso{\wt A}{A \oplus \C}$.
\end{enumerate}

\vskip 8 pt

\noindent\textbf{Kasus 2: aljabar $A$ \emph{tidak} beridentitas.}
\begin{enumerate}
\setcounter{enumi}{8}
    \item Buktikan fakta sederhana berikut.
  \begin{lem} Misalkan $A$ suatu aljabar dan $B$ suatu aljabar bernorma. Jika
  $\phi\colon A \sto B$ adalah homomorfisme aljabar, maka fungsi
  $a \mapsto \norm a := \norm{\phi(a)}$ adalah seminorma pada~$A$. Fungsi tersebut adalah
  norma jika $\phi$ injektif.
  \end{lem}
    \item Ingat bahwa dalam~\ref{0013151} kita mendefinisikan operator $L_a$,
\emph{perkalian kiri oleh $a$}. Sekarang, misalkan
      \[ A^{\sharp} :=
        \{ L_a + \lambda I_A \in \ofml B(A)\colon a \in A \text{ dan } \lambda \in \C\} \]
dan tunjukkan bahwa $A^\sharp$ adalah suatu aljabar bernorma.
    \item Jadikan $A^\sharp$ suatu aljabar-$*\,$ dengan mendefinisikan
      \[ (L_a + \lambda I_A)^* := L_{a^*} + \conj\lambda I_A \]
untuk semua $a \in A$ dan $\lambda \in \C$.
    \item Definisikan
      \[ \phi\colon A \bowtie \C \sto A^\sharp\colon
                   (a,\lambda) \mapsto L_a + \lambda I_A \]
dan verifikasikan bahwa $\phi$ adalah suatu homomorfisme-$*\,$.
    \item Buktikan bahwa $\phi$ injektif.
    \item Gunakan (9) untuk melengkapi $A \bowtie \C$ dengan suatu norma yang
menjadikannya aljabar bernorma beridentitas. Misalkan $\wt A := A \bowtie \C$ dengan
    norma yang ditarik balik oleh $\phi$ dari $A^\sharp$.


    \item Verifikasikan fakta-fakta berikut.
      \begin{enumerate}
        \item Pemetaan $\phi\colon \wt A \sto A^\sharp$ adalah isomorfisme isometrik.
        \item $\ran L$ adalah subaljabar tertutup dari
$\ran \phi = A^\sharp \subseteq \ofml B(A)$.
        \item $I_A \notin \ran L$.
      \end{enumerate}
    \item Buktikan bahwa norma pada $\wt A$ memenuhi syarat-$C^*$.
    \item Buktikan bahwa $\wt A$ adalah suatu aljabar-$C^*$ beridentitas. (Untuk
menunjukkan bahwa $\wt A$ lengkap, kita hanya perlu menunjukkan bahwa $A^\sharp$ lengkap.
Untuk tujuan ini, misalkan $\bigl(\phi(a_n,\lambda_n)\bigr)_{n=1}^\infty$ suatu barisan
Cauchy dalam~$A^\sharp$. Untuk menunjukkan bahwa barisan ini konvergen, cukup ditunjukkan
bahwa barisan itu mempunyai suatu subbarisan konvergen. Dengan menunjukkan bahwa barisan
$(\lambda_n)$ terbatas, kita dapat mengekstrak darinya suatu subbarisan konvergen
$\bigl(\lambda_{n_k}\bigr)$. Buktikan bahwa $(L_{a_{n_k}})$ konvergen.)
    \item Buktikan bahwa barisan
        \begin{equation}\xymatrix{
     0\ar[r] & A\ar[r] & \wt A\ar[r] & \C\ar[r] & 0
        }\end{equation}
eksak terbelah.
    \item Aljabar-$C^*$ $\wt A$ \emph{tidak} isomorfik dengan $A \oplus \C$.
\end{enumerate}  \ns
\end{proof}

\begin{defn}\label{00152131} Aljabar-$C^*$ $\wt A$ yang dikonstruksi dalam proposisi
sebelumnya adalah
 \index{unitalisasi!suatu aljabar-$C^*$}%
\df{unitalisasi} dari~$A$.
\end{defn}

Perhatikan bahwa definisi \emph{spektrum} yang diperluas dalam~\ref{00150014} berlaku untuk
aljabar-$C^*$ karena penambahan identitas merupakan perkara aljabar murni dan sama
bagi aljabar-$C^*$ seperti bagi aljabar Banach umum. Dengan demikian, banyak fakta
terdahulu yang dinyatakan untuk aljabar-$C^*$ beridentitas tetap berlaku. Secara khusus,
untuk rujukan selanjutnya, kita nyatakan kembali butir-butir \ref{00131101},
\ref{00131102}, \ref{00131103}, \ref{00131104}, \ref{001312}, \ref{001406},
dan~\ref{0014062}.

\begin{prop}\label{00152141} Misalkan $a$ suatu unsur normal dari suatu aljabar-$C^*$. Maka
$\norm{a^2} = {\norm a}^2$ dan karena itu $\rho(a)~=~\norm a$.
\end{prop}

\begin{cor}\label{00152142} Jika $A$ suatu aljabar-$C^*$ komutatif, maka
$\norm{a^2} = {\norm a}^2$ dan $\rho(a)~=~\norm a$ untuk setiap $a \in A$.
\end{cor}

\begin{cor}\label{00152143} Pada suatu aljabar-$C^*$ komutatif $A$, transformasi Gelfand
$\Gamma$ adalah isometri; yaitu, $\norm{\Gamma_a}_u = \norm{\wh a}_u = \norm a$ untuk
setiap $a \in A$.
\end{cor}

\begin{cor}\label{00152144} Norma suatu aljabar-$C^*$ bersifat tunggal dalam arti bahwa, jika
diberikan suatu aljabar $A$ dengan involusi, paling banyak ada satu norma yang menjadikan
$A$ suatu aljabar-$C^*$.
\end{cor}

\begin{prop}\label{00152145} Jika $h$ suatu unsur swaadjoin dari suatu aljabar-$C^*$,
 \index{spektrum!unsur swaadjoin}%
 \index{swaadjoin!unsur!spektrum}%
maka $\sigma(h) \subseteq \R$.
\end{prop}

\begin{prop}\label{00152152} Jika $a$ suatu unsur swaadjoin dalam suatu aljabar-$C^*$,
maka transformasi Gelfandnya $\wh a$ bernilai real.
\end{prop}

\begin{prop}\label{00152153} Setiap karakter pada suatu aljabar-$C^*$ $A$ mempertahankan
involusi; dengan demikian, transformasi Gelfand $\sbsb{\Gamma}{\!\!A}$ adalah suatu
homomorfisme-$*$.
\end{prop}

Akibat langsung dari hasil-hasil sebelumnya adalah versi kedua dari
\emph{teorema Gelfand--Naimark}, yang menyatakan bahwa setiap aljabar-$C^*$ komutatif
(isomorfik-$*\,$ secara isometrik dengan) aljabar dari semua fungsi kontinu pada suatu ruang
Hausdorff kompak lokal yang lenyap di tak hingga. Seperti halnya pada versi pertama
teorema ini (lihat~\ref{00142}), ruang Hausdorff kompak lokal yang dimaksud adalah ruang
karakter dari aljabar tersebut.

\begin{thm}[Teorema Gelfand--Naimark II]\label{00152156} Misalkan $A$ suatu aljabar-$C^*$
komutatif.
 \index{Gelfand-Naimark!Teorema II}%
Maka transformasi Gelfand $\Gamma_A\colon a \mapsto \wh a$ adalah suatu
isomorfisme-$*\,$ isometrik dari $A$ pada $\fml C_0(\Delta A)$.
\end{thm}

Dari proposisi berikut, diperoleh bahwa proses unitalisasi bersifat funktorial.

\begin{prop} \label{00152161} Setiap homomorfisme-$*$ $\phi\colon A \sto B$ antara
aljabar-$C^*$ mempunyai suatu perpanjangan tunggal menjadi homomorfisme-$*$ beridentitas
$\wt \phi\colon \wt A \sto \wt B$ antara unitalisasi masing-masing.
\end{prop}

Berbeda dengan keadaan pada aljabar Banach umum, tidak ada perbedaan antara kategori
topologis dan kategori geometris aljabar-$C^*$. Salah satu aspek paling mengagumkan dari
teori aljabar-$C^*$ adalah bahwa homomorfisme-$*$ antara aljabar-aljabar demikian secara
otomatis kontinu---bahkan, kontraktif. Akibatnya, jika dua aljabar-$C^*$ isomorfik-$*$
secara aljabar, maka keduanya isomorfik secara isometrik.

\begin{prop}\label{00152171} Setiap homomorfisme-$*$ antara aljabar-$C^*$ bersifat kontraktif.
\end{prop}

\begin{prop}\label{00152181} Setiap homomorfisme-$*$ injektif antara aljabar-$C^*$ adalah
suatu isometri.
\end{prop}

\begin{prop}\label{00152191} Misalkan $X$ suatu ruang Hausdorff kompak lokal dan
$\wt X = X \cup \{\infty\}$ merupakan kompaktifikasi satu titiknya. Definisikan
  \[ \iota\colon \fml C_0(X) \sto \,\fml C(\wt X)\colon f \mapsto \wt f \]
dengan
 \[ \wt f(x) = \left\{%
  \begin{array}{ll}
     f(x), & \hbox{jika $x \in X$;} \\
     0, & \hbox{jika $x = \infty$.} \\
  \end{array}%
            \right. \]
Definisikan pula $E_\infty$ pada $\fml C(\wt X)$ dengan $E_\infty(g) = g(\infty)$. Maka
barisan
  \[ \vc 0 \to \fml C_0(X) \to^\iota \,\,\fml C(\wt X)
                                        \to^{E_\infty} \C \to \vc 0 \]
eksak.
\end{prop}

Dalam proposisi sebelumnya, kita menyebut $\wt X$ sebagai kompaktifikasi satu titik dari $X$
bahkan jika $X$ sejak awal kompak. Kebanyakan definisi \emph{kompaktifikasi} mensyaratkan
suatu ruang rapat dalam setiap kompaktifikasinya. (Lihat ulasan saya pada awal
bagian 17.3 dari~\cite{Erdman:2007}.) Sebelumnya kita telah menggunakan konvensi bahwa
unitalisasi dari suatu aljabar beridentitas memperoleh suatu identitas perkalian baru.
Demi konsistensi dengan pilihan ini, selanjutnya kita akan menganut konvensi bahwa
kompaktifikasi satu titik dari suatu ruang kompak memperoleh suatu titik tambahan
(terisolasi). (Lihat pula istilah yang kita perkenalkan
dalam~\futurexref{14.3.1}{0038614}.)

 \vskip 2 pt

Dari sudut pandang teorema Gelfand--Naimark~(\ref{00152156}), wawasan mendasar yang
disarankan oleh proposisi berikut adalah bahwa unitalisasi suatu aljabar-$C^*$ komutatif,
dalam arti tertentu, merupakan ``hal yang sama'' dengan kompaktifikasi satu titik dari
suatu ruang Hausdorff kompak lokal.

\begin{prop}\label{00152194} Jika $X$ suatu ruang Hausdorff kompak lokal, maka
aljabar-$C^*$ beridentitas $\bigl(\fml C_0(X)\bigr)^\sim$ dan $\fml C(\wt X)$
isomorfik-$*\,$ secara isometrik.
\end{prop}

\begin{proof} Definisikan
  \[ \theta\colon \bigl(\fml C_0(X)\bigr)^\sim \!\!\sto \fml C(\wt X)\colon
                         (f,\lambda) \mapsto \wt f + \lambda \vc 1_{\wt X} \]
(dengan $\vc 1_{\wt X}$ menyatakan fungsi konstan $1$ pada~$\wt X$). Kemudian tinjau diagram
  \[\xymatrix{
     0\ar[r] & \fml C_0(X)\ar[r]\ar@{=}[d] & \bigl(\fml C_0(X)\bigr)^\sim\ar[r]\ar[d]^\theta
             & \C\ar[r]\ar@{=}[d] & 0 \\
     0\ar[r] & \fml C_0(X)\ar[r]^\iota & \fml C(\wt X)\ar[r]^{E_\infty} & \C\ar[r] & 0
  }\]
Baris atas eksak menurut proposisi~\ref{0015213}, baris bawah eksak menurut
proposisi~\ref{00152191}, dan diagram tersebut jelas komutatif. Pemeriksaan bahwa $\theta$
adalah suatu homomorfisme-$*\,$ bersifat rutin. Oleh karena itu, $\theta$ adalah suatu
isomorfisme-$*\,$ isometrik menurut proposisi~\ref{0015152} dan
korolari~\ref{00152181}.
\end{proof}


















\section{Kuasi-Invers}
\begin{defn} Suatu elemen $b$ dari aljabar $A$ adalah \df{kuasi-invers kiri} bagi $a \in A$ jika
$ba = a + b$. Elemen tersebut adalah \df{kuasi-invers kanan} bagi $a$ jika $ab = a + b$. Jika $b$
sekaligus merupakan kuasi-invers kiri dan kuasi-invers kanan bagi $a$, maka $b$ adalah
 \index{kuasi-invers}%
\df{kuasi-invers} bagi~$a$. Jika $a$ memiliki kuasi-invers, kita lambangkan kuasi-invers itu dengan~$a'$.
\end{defn}

\begin{prop} Jika $b$ adalah kuasi-invers kiri bagi $a$ dalam aljabar $A$ dan $c$ adalah
kuasi-invers kanan bagi $a$ dalam $A$, maka $b = c$.
\end{prop}

\begin{prop} Misalkan $A$ aljabar beridentitas dan $a$, $b \in A$. Maka
 \begin{enumerate}
  \item[(a)] $a'$ ada jika dan hanya jika $(\vc 1 - a)^{-1}$ ada; dan
  \item[(b)] $b^{-1}$ ada jika dan hanya jika $(\vc 1 - b)'$ ada.
 \end{enumerate}
\end{prop}

\begin{proof}[\emph{Petunjuk bukti}] Untuk arah sebaliknya pada (a), tinjau
$a + c - ac$, dengan $c = \vc 1 - (\vc 1 - a)^{-1}$. \ns
\end{proof}

\begin{prop} Suatu elemen $a$ dari aljabar Banach $A$ bersifat kuasi-invertibel jika dan hanya jika
elemen itu bukan merupakan identitas terhadap ideal modular maksimal mana pun dalam~$A$.
\end{prop}

\begin{prop} Misalkan $A$ aljabar Banach dan $a \in A$. Jika $\rho(a) < 1$, maka $a'$ ada
dan $a' = -\sum_{k=1}^\infty a^k$.
\end{prop}

\begin{prop} Misalkan $A$ aljabar Banach dan $a \in A$. Jika $\norm{a} < 1$, maka $a'$ ada
dan
   \[ \frac{\norm{a}}{1 + \norm{a}} \le \norm{a'} \le \frac{\norm{a}}{1 - \norm{a}}. \]
\end{prop}

\begin{prop} Misalkan $A$ aljabar Banach beridentitas dan $a \in A$. Jika $\rho(\vc 1 - a) < 1$, maka
$a$ invertibel dalam~$A$.
\end{prop}

\begin{prop} Misalkan $A$ aljabar Banach dan $a$, $b \in A$. Jika $b'$ ada dan
$\norm{a} < (1 + \norm{b'})^{-1}$, maka $(a + b)'$ ada dan
  \[ \norm{(a + b)' - b'}
           \le \frac{\norm{a}\,(1 + \norm{b'})^2}{1 - \norm{a}\,(1 + \norm{b'})}\,.\]
\end{prop}
 \begin{proof}[\emph{Petunjuk bukti}] Tunjukkan terlebih dahulu bahwa $u = (a - b'a)'$ ada dan bahwa
 $u + b' - ub'$ adalah kuasi-invers kiri bagi~$a~+~b$. \ns
\end{proof}

Bandingkan hasil berikut dengan proposisi~\ref{000644fa} dan~\ref{000646fa}.
\begin{prop} Himpunan $Q_A$ yang terdiri atas
 \index{q@$Q_A$ (elemen kuasi-invertibel)}%
elemen-elemen kuasi-invertibel dari aljabar Banach $A$ adalah himpunan terbuka dalam $A$, dan pemetaan $a \mapsto
a'$ merupakan homeomorfisme dari $Q_A$ ke dirinya sendiri.
\end{prop}

\begin{notn} Jika $a$ dan $b$ adalah elemen dalam suatu aljabar, kita definisikan
 \index{<binop@$a \circ b$ (operasi dalam suatu aljabar)}%
$a \circ b := a + b - ab$.
\end{notn}

\begin{prop} Jika $A$ aljabar, maka $Q_A$ merupakan grup terhadap~$\circ$.
\end{prop}

\begin{prop} Jika $A$ aljabar Banach beridentitas dan $\inv A$ adalah himpunan elemen-elemen
invertibelnya, maka
   \[ \Psi\colon Q_A \sto \inv A\colon  a \mapsto \vc 1 - a \]
merupakan isomorfisme.
\end{prop}

\begin{defn} Jika $A$ aljabar dan $a \in A$, kita definisikan
 \index{spektrum-q}%
 \index{spektrum!q-}%
\df{spektrum-q} dari $a$ dalam $A$ dengan
  \[ \breve{\sigma}_A(a) := \{\lambda \in \C\colon  \lambda \ne 0
                      \text{ dan } \dfrac a\lambda \notin Q_A\}\cup\{0\}. \]
\end{defn}

Dalam aljabar beridentitas, definisi di atas ``hampir'' merupakan definisi yang biasa.
\begin{prop} Misalkan $A$ aljabar beridentitas dan $a \in A$. Maka untuk setiap $\lambda \ne 0$, berlaku
$\lambda \in \sigma(a)$ jika dan hanya jika $\lambda \in \breve{\sigma}(a)$.
\end{prop}

\begin{prop} Jika aljabar $A$ tak beridentitas dan $a \in A$, maka $\breve{\sigma}_A(a) =
\breve{\sigma}_{\wt A}(a)$, dengan $\wt A$ adalah unitalisasi dari~$A$.
\end{prop}























\section{Elemen Positif dalam Aljabar-$C^*$}
\begin{defn} Suatu elemen swaadjoin $a$ dari aljabar-$C^*$ $A$ disebut
 \index{positif!elemen dari suatu aljabar-$C^*$}%
\df{positif} jika $\sigma(a) \subseteq [0,\infty)$. Dalam hal ini kita tulis $a \ge \vc 0$.
Kita lambangkan himpunan semua elemen positif dari $A$
 \index{<plus@$A^+$ (kerucut positif)}%
dengan~$A^+$. Himpunan ini adalah
 \index{positif!kerucut}%
 \index{kerucut!positif}%
\df{kerucut positif} dari~$A$. Untuk setiap himpunan bagian $B$ dari $A$, misalkan $B^+ = B \cap A^+$. Kita akan
menggunakan kerucut positif untuk menginduksi suatu urutan parsial pada~$A$: kita tulis
 \index{<relation@$a \le b$ (relasi urutan dalam suatu aljabar-$C^*$)}%
$a \le b$ apabila $b - a \in A^+$.
\end{defn}

\begin{defn} Misalkan $\le$ suatu relasi pada himpunan tak kosong $S$. Jika relasi
 \index{<relation@$\le$ (notasi untuk suatu urutan parsial)}%
 $\le$ bersifat refleksif dan transitif, relasi tersebut adalah suatu
 \index{praurutan}%
 \index{urutan!pra-}%
\df{praurutan}. Jika $\le$ adalah praurutan dan juga antisimetris, relasi tersebut adalah suatu
 \index{parsial!urutan}%
 \index{urutan!parsial}%
\df{urutan parsial}.

Suatu urutan parsial $\le$ pada ruang vektor riil $V$
 \index{kompatibilitas!urutan dengan operasi}%
\df{kompatibel} dengan (atau
 \index{menghormati operasi}%
\df{menghormati}) operasi-operasi (penjumlahan dan perkalian skalar) pada~$V$ jika untuk semua $x$,
$y$, $z \in V$
 \begin{enumerate}
  \item[(a)] $x \le y$ mengakibatkan $x+z \le y+z$, dan
  \item[(b)] $x \le y$, $\alpha \ge 0$ mengakibatkan $\alpha x \le \alpha y$.
 \end{enumerate}
Suatu ruang vektor riil yang dilengkapi dengan urutan parsial yang kompatibel dengan operasi-operasi
ruang vektor disebut
 \index{terurut!ruang vektor}%
 \index{vektor!ruang!terurut}%
\df{ruang vektor terurut}.
\end{defn}

\begin{defn} Misalkan $V$ ruang vektor. Suatu himpunan bagian $C$ dari $V$ adalah
 \index{kerucut}%
\df{kerucut} dalam $V$ jika $\alpha C \subseteq C$ untuk setiap $\alpha \ge 0$. Suatu kerucut $C$ dalam $V$
disebut
 \index{proper!kerucut}%
 \index{kerucut!proper}%
\df{proper} jika $C \cap (-C) = \{\vc 0\}$.
\end{defn}

\begin{prop}\label{0018006} Suatu kerucut $C$ dalam ruang vektor bersifat konveks jika dan hanya jika $C + C
\subseteq C$.
\end{prop}

\begin{exam}\label{0018007} Jika $V$ ruang vektor terurut, maka himpunan
 \index{<dualp@$V^+$ (kerucut positif dalam suatu ruang vektor terurut)}%
   \[ V^+ := \{x \in V \colon x \ge \vc 0\} \]
adalah kerucut konveks proper dalam~$V$. Himpunan ini adalah
 \index{positif!kerucut!dalam suatu ruang vektor terurut}%
 \index{kerucut!positif}%
\df{kerucut positif} dari~$V$, dan anggota-anggotanya adalah
 \index{positif!elemen (dari suatu ruang vektor terurut)}%
\df{elemen positif} dari~$V$.
\end{exam}

\begin{prop}\label{0018008} Misalkan $V$ ruang vektor riil dan $C$ kerucut konveks proper
dalam~$V$. Definisikan $x \le y$ jika $y - x \in C$. Maka relasi $\le$ adalah urutan parsial
pada $V$ dan kompatibel dengan operasi-operasi ruang vektor pada~$V$. Relasi ini adalah
\df{urutan parsial
 \index{terinduksi!urutan parsial}%
 \index{parsial!urutan!terinduksi oleh kerucut konveks proper}%
yang diinduksi oleh} kerucut~$C$. Kerucut positif $V^+$ dari ruang vektor terurut yang dihasilkan
tepat sama dengan $C$.
\end{prop}

\begin{prop} Jika $a$ adalah elemen swaadjoin dari suatu aljabar-$C^*$ beridentitas dan $t \in \R$, maka
 \begin{enumerate}
   \item[(i)] $a \ge \vc 0$ apabila $\norm{a - t\vc 1} \le t$; dan
   \item[(ii)] $\norm{a - t\vc 1} \le t$ apabila $\norm a \le t$ dan $a \ge \vc 0$.
  \end{enumerate}
\end{prop}

\begin{exam} Kerucut positif dari suatu aljabar-$C^*$ $A$ adalah kerucut konveks proper tertutup
dalam ruang vektor riil~$\ofml H(A)$.
\end{exam}

\begin{prop}\label{0018015} Jika $a$ dan $b$ adalah elemen positif dari suatu aljabar-$C^*$
dan $ab = ba$, maka $ab$ positif.
\end{prop}

\begin{prop}\label{001804} Setiap elemen positif dari suatu aljabar-$C^*$ $A$ memiliki akar pangkat
$n$ positif yang tunggal ($n \in \N$). Artinya, jika $a \in A^+$, maka terdapat tepat satu
$b \in A^+$ sedemikian sehingga $b^n = a$.
\end{prop}

\begin{proof} \emph{Petunjuk.} Bagian eksistensi merupakan penerapan sederhana kalkulus fungsional
$C^*$ (yaitu, \emph{teorema spektral abstrak}~\ref{00144}). Elemen $b$
yang diberikan oleh kalkulus fungsional bersifat positif dalam aljabar $C^*(\vc 1,a)$. Jelaskan mengapa
elemen itu juga positif dalam~$A$. Argumen keunikan memerlukan perhatian yang saksama. \ns
\end{proof}

\begin{thm}[Dekomposisi Jordan]\label{001807} Jika $c$ adalah elemen swaadjoin dari suatu
 \index{dekomposisi Jordan}%
 \index{dekomposisi!Jordan}%
aljabar-$C^*$ $A$, maka terdapat elemen-elemen positif
 \index{<positive@$c^+$ (bagian positif dari $c$)}%
 \index{positif!bagian!dari suatu elemen swaadjoin}%
 \index{<positive@$c^-$ (bagian negatif dari $c$)}%
 \index{negatif!bagian!dari suatu elemen swaadjoin}%
$c^+$ dan $c^-$ dalam $A$ sedemikian sehingga $c = c^+ - c^-$ dan $c^+c^- = \vc 0$.
\end{thm}

\begin{lem} Jika $c$ adalah elemen dari suatu aljabar-$C^*$ sedemikian sehingga $-c^*c \ge \vc 0$, maka $c
= \vc 0$.
\end{lem}

\begin{prop} Jika $a$ adalah elemen dari suatu aljabar-$C^*$, maka $a^*a \ge \vc 0$.
\end{prop}

\begin{prop}\label{00180731} Jika $c$ adalah elemen dari suatu aljabar-$C^*$ $A$, maka pernyataan-pernyataan
berikut ekuivalen:
 \begin{enumerate}
    \item[(i)] $c \ge \vc 0$;
    \item[(ii)] terdapat $b \ge \vc 0$ sedemikian sehingga $c = b^2$; dan
    \item[(iii)] terdapat $a \in A$ sedemikian sehingga $c = a^*a$,
 \end{enumerate}
\end{prop}

\begin{exam}\label{00180741} Jika $T$ operator pada ruang Hilbert $H$, maka $T$ merupakan
elemen positif dari aljabar-$C^*$ $\ofml B(H)$ jika dan hanya jika $\langle Tx,x \rangle
\ge 0$ untuk setiap $x \in H$.
\end{exam}

\begin{proof}[\emph{Petunjuk bukti}] Mudah untuk menunjukkan bahwa jika $T \ge \vc 0$ dalam $\ofml B(H)$, maka
$\langle Tx,x \rangle \ge 0$ untuk setiap $x \in H$: gunakan proposisi~\ref{00180731} untuk
menuliskan $T$ sebagai $S^*S$ bagi suatu operator~$S$.

Sebaliknya, andaikan $\langle Tx,x \rangle \ge 0$ untuk setiap $x \in H$. Mudah
dilihat bahwa ini mengakibatkan $T$ swaadjoin. Gunakan \emph{teorema dekomposisi
Jordan}~\ref{001807} untuk menuliskan $T$ sebagai $T^+ - T^-$. Untuk sembarang $u \in H$, misalkan $x = T^-u$
dan periksa bahwa $0 \le \langle Tx,x \rangle = -\langle (T^-)^3u,u \rangle$. Sekarang $(T^-)^3$
adalah elemen positif dari $\ofml B(H)$. (Mengapa?) Simpulkan bahwa $(T^-)^3 = \vc 0$ dan
karena itu $T^- = \vc 0$. (Untuk rincian tambahan, lihat~\cite{DoranB:1986}, halaman~37.) \ns
\end{proof}

\begin{defn} Untuk sembarang elemen $a$ dari suatu aljabar-$C^*$, kita definisikan
 \index{mutlak!nilai!dalam suatu aljabar-$C^*$}%
 \index{<unop@$\abs a$ (nilai mutlak dalam suatu aljabar-$C^*$)}%
$\abs a$ sebagai $\sqrt{a^*a}$.
\end{defn}

\begin{prop} Jika $a$ adalah elemen swaadjoin dari suatu aljabar-$C^*$, maka
  \[ \abs a = a^+ + a^-\,. \]
\end{prop}

\begin{exam} Nilai mutlak dalam suatu aljabar-$C^*$ tidak harus
subaditif; yaitu, $\abs{a + b}$ tidak harus lebih kecil daripada $\abs{a}
+ \abs{b}$. Sebagai contoh, dalam $\M 2\C$, ambil $a = \begin{bmatrix} 1 & 1 \\
1 & 1 \end{bmatrix}$ dan $b = \begin{bmatrix} 0 & 0 \\ 0 & -2
\end{bmatrix}$. Maka $\abs{a} = a$, $\abs{b} = \begin{bmatrix} 0 &
0 \\ 0 & 2 \end{bmatrix}$, dan $\abs{a + b} = \begin{bmatrix} \sqrt{2} & 0 \\ 0 &
\sqrt{2}
\end{bmatrix}$. Jika $\abs{a + b} - \abs{a} - \abs{b}$ positif, maka, berdasarkan
contoh~\ref{00180741}, $\langle\,(\abs{a + b} - \abs{a} - \abs{b})\,x\,,\,x\,\rangle$
akan positif untuk setiap vektor $x \in \C^2$. Namun, hal ini tidak benar untuk $x = (1,0)$.
\end{exam}

\begin{prop}\label{0018201} Jika $\phi\colon A \sto B$ adalah homomorfisme-$*\,$ antara
aljabar-$C^*$, maka $\phi(a) \in B^+$ apabila $a \in A^+$. Jika $\phi$ adalah
isomorfisme-$*\,$, maka $\phi(a) \in B^+$ jika dan hanya jika $a \in A^+$.
\end{prop}

\begin{prop} Misalkan $a$ elemen swaadjoin dari suatu aljabar-$C^*$ $A$ dan $f$ fungsi kontinu
bernilai kompleks pada spektrum~$a$. Maka $f \ge \vc 0$ dalam $\fml
C(\sigma(a))$ jika dan hanya jika $f(a) \ge \vc 0$ dalam~$A$.
\end{prop}

\begin{prop}\label{0018203} Jika $a$ adalah elemen swaadjoin dari suatu aljabar-$C^*$ $A$, maka
$\norm a\,\vc 1_{\wt A} \pm a \ge \vc 0$ dalam~$\wt A$.
\end{prop}

\begin{prop}\label{0018204} Jika $a$ dan $b$ adalah elemen swaadjoin dari suatu
aljabar-$C^*$~$A$ dan $a \le b$, maka $x^*ax \le x^*bx$ untuk setiap $x \in A$.
\end{prop}

\begin{prop} Jika $a$ dan $b$ adalah elemen dari suatu aljabar-$C^*$ dengan $0 \le a \le b$, maka
$\norm a \le \norm b$.
\end{prop}

\begin{prop}\label{0018301} Misalkan $A$ aljabar-$C^*$ beridentitas dan $c \in A^+$. Maka
$c$ invertibel jika dan hanya jika $c \ge \epsilon\vc 1$ untuk suatu $\epsilon > 0$.
\end{prop}

\begin{prop} Misalkan $A$ aljabar-$C^*$ beridentitas dan $c \in A$. Jika $c \ge \vc 1$, maka $c$
invertibel dan $\vc 0 \le c^{-1} \le \vc 1$.
\end{prop}

\begin{prop} Jika $a$ adalah elemen positif yang invertibel dalam suatu aljabar-$C^*$ beridentitas, maka $a^{-1}$
positif.
\end{prop}

Selanjutnya kita tunjukkan bahwa notasi $a^{\ssst{-\frac12}}$ tidak ambigu.

\begin{prop} Misalkan $a \in A^+$, dengan $A$ aljabar-$C^*$ beridentitas. Jika $a$ invertibel, maka
$a^{\ssst{\frac12}}$ juga invertibel dan $\bigl(a^{\ssst{\frac12}}\bigr)^{-1} =
\bigl(a^{-1}\bigr)^{\ssst{\frac12}}$.
\end{prop}

\begin{cor}\label{00183313} Jika $a$ adalah elemen invertibel dalam suatu aljabar-$C^*$ beridentitas,
 \index{invers!dari suatu nilai mutlak}%
maka~$\abs a$ juga invertibel.
\end{cor}

\begin{prop}\label{0018401} Misalkan $a$ dan $b$ adalah elemen dari suatu aljabar-$C^*$. Jika
$\vc 0 \le a \le b$ dan $a$ invertibel, maka $b$ invertibel dan $b^{-1} \le
a^{-1}$.
\end{prop}

\begin{prop} Jika $a$ dan $b$ adalah elemen dari suatu aljabar-$C^*$ dan $\vc 0 \le a \le b$,
maka $\sqrt a \le \sqrt b$.
\end{prop}

\begin{exam} Misalkan $a$ dan $b$ elemen dari suatu aljabar-$C^*$ dengan $\vc 0 \le a \le b$.
\emph{Tidak} selalu berlaku bahwa $a^2 \le b^2$.
\end{exam}

\begin{proof}[\emph{Petunjuk bukti}] Dalam aljabar-$C^*$ $\mathbf M_2$, misalkan $a =
\begin{bmatrix} 1 & 0 \\ 0 & 0 \end{bmatrix}$ dan
$b = a + \tfrac12 \begin{bmatrix} 1 & 1 \\ 1 & 1 \end{bmatrix}$. \ns
\end{proof}
























\section{Identitas Aproksimatif}
\begin{defn} Suatu
 \index{identitas!aproksimatif}%
 \index{aproksimatif!identitas}%
\df{identitas aproksimatif} (atau
 \index{unit!aproksimatif}%
 \index{aproksimatif!unit}%
\df{unit aproksimatif}) dalam suatu aljabar-$C^*$ $A$ adalah jaring menaik
${(e_\lambda)}_{\lambda \in \Lambda}$ yang terdiri atas elemen-elemen positif dari $A$ sedemikian sehingga
$\norm{e_\lambda} \le 1$ dan $ae_\lambda \sto a$ (secara ekuivalen, $e_\lambda\, a \sto a$)
untuk setiap $a \in A$. Jika jaring tersebut ternyata merupakan barisan, kita memperoleh
 \index{sekuensial!identitas aproksimatif}%
 \index{aproksimatif!identitas!sekuensial}%
\df{identitas aproksimatif sekuensial}. Dalam literatur, berhati-hatilah terhadap definisi yang
beragam: banyak penulis menghilangkan syarat bahwa jaring tersebut menaik dan/atau bahwa jaring itu
terbatas.
\end{defn}

\begin{exam} Misalkan $A = \fml C_0(\R)$. Untuk setiap $n \in \N$, misalkan $U_n = (-n,n)$ dan
$e_n\colon \R \sto [0,1]$ suatu fungsi dalam $A$ yang dukungannya termuat dalam $U_{n+1}$
dan sedemikian sehingga $e_n(x) = 1$ untuk setiap $x \in U_n$. Maka $(e_n)$ adalah identitas
aproksimatif (sekuensial) bagi~$A$.
\end{exam}

\begin{prop} Jika $A$ adalah aljabar-$C^*$, maka himpunan
 \index{lambda@$\Lambda$ (anggota $A^+$ dalam bola satuan terbuka)}%
   \[ \Lambda := \{a \in A^+\colon \norm a < 1\} \]
merupakan himpunan terarah (terhadap urutan yang diwarisinya dari $\ofml H(A)$).
\end{prop}

\begin{proof}[\emph{Petunjuk bukti}] Periksa bahwa fungsi
  \[ f\colon [0,1) \sto \R^+\colon t \mapsto (1 - t)^{-1} - 1 = \frac t{1 - t} \]
menginduksi isomorfisme urutan $f\colon \Lambda \sto A^+$. (Suatu
 \index{urutan!isomorfisme}%
 \index{isomorfisme!urutan}%
\df{isomorfisme urutan} antara himpunan-himpunan terurut parsial adalah bijeksi yang mempertahankan urutan,
yang inversnya juga mempertahankan urutan.) Bukti yang saksama untuk hasil ini melibatkan pemeriksaan
terhadap cukup banyak rincian. \ns
\end{proof}

\begin{prop} Jika $A$ adalah aljabar-$C^*$, maka himpunan
  \[ \Lambda := \{a \in A^+\colon \norm a < 1\} \]
merupakan identitas aproksimatif bagi~$A$.
\end{prop}

\begin{cor}\label{00190161} Setiap aljabar-$C^*$ $A$ memiliki identitas aproksimatif.
 \index{aproksimatif!identitas!keberadaan}%
Jika $A$ separabel, maka $A$ memiliki identitas aproksimatif sekuensial.
\end{cor}

\begin{prop}\label{0019017} Setiap ideal aljabar (dua sisi) tertutup dalam suatu aljabar-$C^*$
bersifat swaadjoin.
\end{prop}

\begin{prop}\label{001902} Jika $J$ ideal dalam suatu aljabar-$C^*$ $A$, maka
 \index{hasil bagi!aljabar-$C^*$}%
$A/J$ adalah aljabar-$C^*$.
\end{prop}

\begin{proof} Lihat \cite{HigsonR:2000}, teorema 1.7.4, atau \cite{Davidson:1996}, halaman
13--14, atau~\cite{Fillmore:1996}, teorema 2.5.4. \ns
\end{proof}

\begin{exam} Jika $J$ ideal dalam suatu aljabar-$C^*$ $A$, maka barisan
  \[ \vc 0 \to J \to A \to^\pi A/J \to \vc 0 \]
adalah barisan eksak pendek.
\end{exam}

\begin{prop}\label{0019023} Jangkauan suatu homomorfisme-$*$ antara aljabar-$C^*$ bersifat tertutup
(dan karena itu merupakan aljabar-$C^*$).
\end{prop}

\begin{defn} Suatu subaljabar-$C^*$ $B$ dari aljabar-$C^*$ $A$ disebut
 \index{herediter}%
\df{herediter} jika $a \in B$ setiap kali $a \in A$, $b \in B$, dan $\vc 0 \le a \le b$.
\end{defn}

\begin{prop}\label{001915} Andaikan $x^*x \le a$ dalam suatu aljabar-$C^*$ $A$. Maka terdapat
$b \in A$ sedemikian sehingga $x = ba^{\frac14}$ dan $\norm b \le {\norm a}^\frac14$.
\end{prop}

\begin{proof} Lihat \cite{Davidson:1996}, halaman 13. \ns
\end{proof}

\begin{prop} Andaikan $J$ ideal dalam suatu aljabar-$C^*$ $A$, $j \in J^+$, dan
$a^*a \le j$. Maka $a \in J$. Jadi, ideal-ideal dalam aljabar-$C^*$ bersifat herediter.
\end{prop}

\begin{thm}\label{001922} Misalkan $A$ dan $B$ aljabar-$C^*$ dan $J$ ideal dalam~$A$.\newline
Jika $\phi$ \mbox{homomorfisme-$*\,$} dari $A$ ke $B$ dan $\ker\phi \supseteq J$, maka terdapat
tepat satu \mbox{homomorfisme-$*\,$} $\widetilde\phi\colon A/J \sto B$ yang membuat
diagram berikut komutatif.
 \[ \xymatrix@+30pt{A \ar[d]_*+{\pi} \ar[dr]^*+{\phi} & \\
            A/J \ar[r]_*+{\widetilde\phi} & B  } \]
Selain itu, $\widetilde\phi$ injektif jika dan hanya jika $\ker\phi = J$; dan
$\widetilde\phi$ surjektif jika dan hanya jika $\phi$ surjektif.
\end{thm}

\begin{cor}\label{001923} Jika $\vc 0 \to A \to^\phi E \to B \to \vc 0$ adalah barisan eksak
pendek dari aljabar-$C^*$, maka $E/\ran\phi$ dan $B$ isomorfik-$*\,$ secara isometrik.
\end{cor}

\begin{cor}\label{001924} Setiap aljabar-$C^*$ $A$ memiliki kodimensi satu dalam
unitalisasinya~$\wt A$; yaitu, $\dim \wt A/A = 1$.
\end{cor}

\begin{prop} Misalkan $A$ aljabar-$C^*$, $B$ subaljabar-$C^*$ dari $A$, dan $J$
ideal dalam~$A$. Maka
   \[  B/(B \cap J) \cong (B + J)/J\,. \]
\end{prop}

Misalkan $B$ subaljabar beridentitas dari suatu aljabar sembarang $A$. Jelas bahwa jika suatu
elemen $b \in B$ invertibel dalam $B$, maka elemen itu juga invertibel dalam~$A$. Ternyata
konversnya benar dalam aljabar-$C^*$: jika $b$ invertibel dalam $A$, maka inversnya
terletak dalam~$B$. Hal ini biasanya dinyatakan dengan mengatakan bahwa \emph{setiap subaljabar-$C^*$
beridentitas dari suatu aljabar-$C^*$ bersifat}
 \index{invers!tertutup}%
 \index{tertutup!invers}%
\df{tertutup terhadap invers}.

\begin{prop} Misalkan $B$ subaljabar-$C^*$ beridentitas dari suatu aljabar-$C^*$ $A$. Jika $b \in \inv(A)$,
maka $b^{-1} \in B$.
\end{prop}

\begin{cor}\label{0019283} Misalkan $B$ subaljabar-$C^*$ beridentitas dari suatu aljabar-$C^*$ $A$ dan
$b \in B$. Maka
   \[ \sigma_B(b) = \sigma_A(b)\,. \]
\end{cor}

\begin{cor}\label{0019286} Misalkan $\phi\colon A \sto B$ adalah monomorfisme-$C^*$ beridentitas dari
suatu aljabar-$C^*$ $A$ dan $a \in A$. Maka
   \[ \sigma(a) = \sigma(\phi(a))\,. \]
\end{cor}



\endinput
